\chapter[1984 --- Jerne et al.]{1984: Niels K. Jerne, Georges J.F. K\"ohler, C\'esar Milstein}
\label{ch:1984_Jerne}

\section*{Main Contribution}

\textit{Developed the technique for producing monoclonal antibodies, enabling the creation of highly specific diagnostic and therapeutic tools.}

\section{Official Citation}

for theories concerning the specificity in development and control of the immune system and the discovery of the principle for production of monoclonal antibodies

The 1984 Nobel Prize in Physiology or Medicine was awarded jointly to Niels K. Jerne, Georges J.F. K\"ohler, and C\'esar Milstein for their contributions to immunology: Jerne for his theoretical frameworks concerning the specificity and regulation of the immune system, and K\"ohler and Milstein for their development of the technique for producing monoclonal antibodies. Together, their work transformed immunology from a descriptive science into a molecular discipline and provided tools---monoclonal antibodies---that have revolutionized medicine, research, and biotechnology.

\section{Background and Historical Context}

The immune system is one of the most complex and elegant biological systems in the body. Its primary function is to distinguish between ``self'' (the body's own tissues) and ``non-self'' (foreign invaders such as bacteria, viruses, and parasites) and to mount targeted responses against the latter while sparing the former. The adaptive immune system---the branch of the immune system that produces antibodies and T cell receptors---is characterized by its notable diversity and specificity: it can recognize and respond to millions of different foreign molecules (antigens), and it can remember previous encounters with these antigens, mounting a faster and more effective response upon re-exposure.

The central question of immunology in the twentieth century was: how does the immune system generate this notable diversity of antibody molecules? By the 1960s, it was known that antibodies are Y-shaped proteins produced by B lymphocytes, and that each antibody molecule has a specific binding site (paratope) that recognizes a specific antigenic determinant (epitope). However, the mechanism by which the immune system could generate millions of different antibody specificities---each encoded by a different gene---was unknown. The prevailing ``one gene, one polypeptide'' hypothesis of the time could not account for the vast diversity of antibodies, because the entire human genome does not contain enough genes to encode millions of different antibody molecules.

Niels Kaj Jerne was born in London in 1911 to Danish parents, and he grew up in Denmark. He trained in medicine at the University of Copenhagen but abandoned clinical practice to pursue a career in research. Jerne's early work was on the theoretical aspects of antibody formation, and he became one of an influential theoretical immunologists of the twentieth century.

Georges Jean Fran\c{c}ois K\"ohler was born in M\"unster, Germany, in 1946, and trained in biology at the University of Freiburg. He joined the Laboratory of Molecular Biology in Cambridge, England, in 1974, where he collaborated with C\'esar Milstein.

C\'esar Milstein was born in Bah\'ia Blanca, Argentina, in 1927, and trained in chemistry and biochemistry at the University of Buenos Aires. He emigrated to England in 1961 and joined the Medical Research Council (MRC) Laboratory of Molecular Biology in Cambridge, where he conducted his Nobel Prize-winning work on monoclonal antibodies.

\section{The Discovery/Contribution}

\subsection{Niels K. Jerne and the Theoretical Framework of Immunology}

Jerne's contributions to immunology were primarily theoretical, but they were of fundamental importance. In the 1950s and 1960s, Jerne proposed several significant ideas that transformed our understanding of how the immune system works.

Jerne's first major contribution was the ``natural selection theory of antibody formation,'' published in 1955. At the time, the dominant theory of antibody formation was the ``instruction theory,'' which held that antigens served as templates that instructed the immune system to produce complementary antibodies. Jerne proposed an alternative: that the immune system produces a diverse repertoire of antibody molecules spontaneously, in the absence of any antigen, and that antigens merely select those antibody molecules that happen to bind to them. This theory---which was conceptually similar to Darwin's theory of natural selection---was a radical departure from the instruction theory, and it was initially met with skepticism.

However, Jerne's natural selection theory was largely vindicated by the subsequent discovery of the genetic mechanisms of antibody diversity (see below, and Chapter~1987). It is now well established that the immune system generates a vast repertoire of antibody specificities through genetic mechanisms---including somatic recombination, somatic hypermutation, and gene conversion---that operate independently of antigen stimulation. Antigens then select those B lymphocytes whose antibodies happen to bind to them, and these selected B lymphocytes proliferate and differentiate into antibody-secreting plasma cells.

Jerne's second major contribution was the ``network theory of the immune system,'' published in 1974. In this theory, Jerne proposed that the immune system is regulated by a complex network of interactions between antibodies and anti-antibodies (antibodies that recognize the antigen-binding sites of other antibodies). According to the network theory, each antibody molecule has two faces: an idiotypic face (the antigen-binding site) and an isotypic face (the constant region). The idiotypic face can be recognized by other antibodies (anti-idiotypic antibodies), which in turn can be recognized by anti-anti-idiotypic antibodies, and so on. This network of interactions, Jerne argued, provides a mechanism for the regulation of the immune response---maintaining the immune system in a state of dynamic equilibrium and preventing excessive or inappropriate immune activation.

The network theory was controversial and has been modified over the years, but it introduced the important concept that the immune system is not simply a collection of independent antibodies and cells, but rather a self-regulating network in which the components interact with and influence each other. This concept has influenced our understanding of autoimmune disease, immunological tolerance, and the design of vaccines.

\subsection{Georges K\"ohler and C\'esar Milstein and the Production of Monoclonal Antibodies}

K\"ohler and Milstein's development of the hybridoma technique for producing monoclonal antibodies was one of the major methodological innovations in the history of biology. The technique, first described in a landmark paper published in \emph{Nature} in 1975, made it possible to produce unlimited quantities of antibodies of a single specificity---monoclonal antibodies---by fusing antibody-producing B lymphocytes with immortal myeloma (cancer) cells.

The problem that K\"ohler and Milstein set out to solve was this: when an animal is immunized with an antigen, its immune system produces a polyclonal response---that is, many different B lymphocyte clones are activated, each producing an antibody with a different specificity. The resulting antiserum is a mixture of many different antibodies, only some of which are specific for the antigen of interest. To obtain a pure, specific antibody, it is necessary to isolate a single B lymphocyte clone and to grow it in culture so that it produces a homogeneous population of identical antibody molecules.

However, normal B lymphocytes cannot be grown indefinitely in culture---they have a limited lifespan and die after a few cell divisions. K\"ohler and Milstein's insight was to fuse B lymphocytes with myeloma cells---cancerous B lymphocytes that can grow indefinitely in culture---to create hybrid cells (hybridomas) that combine the antibody-producing ability of the B lymphocyte with the immortality of the myeloma cell. By fusing B lymphocytes from an immunized mouse with myeloma cells, selecting the hybrid cells, and screening them for the production of the desired antibody, K\"ohler and Milstein were able to isolate hybridoma clones that produced a single, specific antibody in unlimited quantities.

The practical significance of monoclonal antibodies was immediately recognized. Within a few years of the publication of K\"ohler and Milstein's paper, monoclonal antibodies were being used in research laboratories around the world for a wide range of applications, including the detection and purification of proteins, the diagnosis of diseases, and the study of cellular and molecular processes. The technique was relatively simple, inexpensive, and reproducible, and it democratized the production of specific antibodies---previously a laborious and unreliable process---by making it available to any laboratory with access to a tissue culture facility.

\section{Impact on Modern Medicine}

The impact of monoclonal antibodies on modern medicine has been transformative. Monoclonal antibodies are now among a major classes of therapeutic agents in medicine, with more than 100 monoclonal antibody drugs approved for clinical use and many more in development.

In oncology, monoclonal antibodies have revolutionized the treatment of cancer. Rituximab (Rituxan), a monoclonal antibody that targets the CD20 antigen on B lymphocytes, is used to treat non-Hodgkin lymphoma, chronic lymphocytic leukemia, and autoimmune diseases. Trastuzumab (Herceptin), a monoclonal antibody that targets the HER2 receptor, is used to treat HER2-positive breast cancer. Bevacizumab (Avastin), a monoclonal antibody that targets vascular endothelial growth factor (VEGF), is used to inhibit the growth of new blood vessels in tumors. Cetuximab (Erbitux), a monoclonal antibody that targets the epidermal growth factor receptor (EGFR), is used to treat colorectal cancer and head and neck cancer. These monoclonal antibody drugs have significantly improved survival and quality of life for many cancer patients.

In autoimmune disease, monoclonal antibodies have provided new options for patients with rheumatoid arthritis, Crohn's disease, psoriasis, and other conditions. Adalimumab (Humira), a monoclonal antibody that targets tumor necrosis factor alpha (TNF-$\alpha$), is one of the most widely prescribed drugs in the world and is used to treat a range of autoimmune conditions. Infliximab (Remicade), etanercept (Enbrel), and other TNF inhibitors have similarly transformed the management of autoimmune diseases.

In infectious disease, monoclonal antibodies have been used for the prevention and treatment of respiratory syncytial virus (RSV) infection in infants (palivizumab/Synagis), for the treatment of anthrax exposure (raxibacumab), and for the prevention of SARS-CoV-2 infection (bamlanivimab, casirivimab/imdevimab). The development of monoclonal antibody therapies for COVID-19 demonstrated the speed and versatility of the monoclonal antibody platform in responding to emerging infectious diseases.

Monoclonal antibodies are also widely used in diagnostic medicine. The pregnancy test, the HIV test, and many other diagnostic tests are based on the use of monoclonal antibodies to detect specific antigens or antibodies in blood and other body fluids. Monoclonal antibodies are also used in immunohistochemistry to identify specific cell types and proteins in tissue sections, and in flow cytometry to characterize and sort cells based on their surface markers.

The success of monoclonal antibody therapy has led to the development of new generations of antibody-based therapeutics, including bispecific antibodies (which can bind two different antigens simultaneously), antibody-drug conjugates (which combine the targeting ability of antibodies with the cytotoxicity of chemotherapy drugs), and nanobodies (small, single-domain antibodies derived from camelid antibodies). These next-generation antibody therapeutics are expanding the range of diseases that can be treated with antibody-based approaches.

\section{Legacy}

Niels Jerne continued to develop his theories of immune regulation at the Basel Institute for Immunology in Switzerland, where he served as director from 1980 to 1985. He died in 1993, but his theoretical frameworks continue to influence immunological thinking.

Georges K\"ohler continued his work at the Basel Institute for Immunology and later at the German Cancer Research Center in Heidelberg, where he made further contributions to monoclonal antibody technology. He died unexpectedly in 1995 at the age of 48, but his work continues to save lives around the world.

C\'esar Milstein continued his work at the MRC Laboratory of Molecular Biology in Cambridge until his retirement. He continued to make contributions to antibody engineering and to the development of new antibody-based therapeutics. He died in 2002, but his work on monoclonal antibodies remains one of the most consequential discoveries in the history of biology.

The work of the three laureates of 1984 transformed immunology from a descriptive science into a molecular and therapeutic discipline. Jerne's theoretical insights revealed the elegance and complexity of the immune system, while K\"ohler and Milstein's monoclonal antibody technology provided tools that have revolutionized medicine, research, and biotechnology. The monoclonal antibody is, by any measure, one of the major inventions in the history of medicine---a discovery that has saved and improved countless lives and that continues to generate new therapeutic possibilities.


\section{Modern Developments}

Jerne, Köhler, and Milstein's monoclonal antibody work has led to modern antibody therapeutics. Over 100 monoclonal antibodies are approved for clinical use, treating cancers, autoimmune diseases, and infectious diseases. Antibody-drug conjugates (ADCs) deliver chemotherapy specifically to cancer cells. Bispecific antibodies engage immune cells to attack tumors. Phage display and yeast display enable rapid antibody discovery. Understanding of antibody engineering has led to humanized and fully human antibodies.


\section{Bibliography}

\begin{thebibliography}{9}
\bibitem{nobel-1984}
Nobel Prize Outreach, ``The Nobel Prize in Physiology or Medicine 1984,''\emph{NobelPrize.org}.\\
\url{https://www.nobelprize.org/prizes/medicine/1984/summary/}

\bibitem{lecture-1984}
Niels K. Jerne, ``Nobel Lecture,''\emph{NobelPrize.org}. Winner-authored primary source; downloadable from the official Nobel Prize archive.\\
\url{https://www.nobelprize.org/prizes/medicine/1984/jerne/lecture/}
\end{thebibliography}
